\documentclass[aspectratio=43,11pt]{beamer}
\usepackage[utf8]{inputenc}
\usepackage[russian]{babel}
\usepackage{graphicx}
\usepackage{color}
\usepackage{listings}
\usepackage{xcolor}
\usepackage{hyperref}
\usepackage{geometry}

% Установка светлой темы
\usetheme{default}
\usecolortheme{default}

% Кастомизация цветов
\setbeamercolor{background canvas}{bg=white}
\setbeamercolor{normal text}{fg=black}
\setbeamercolor{frametitle}{fg=darkblue, bg=white}
\setbeamercolor{title}{fg=darkblue}
\setbeamercolor{subtitle}{fg=gray}
\setbeamercolor{itemize item}{fg=darkblue}
\setbeamercolor{enumerate item}{fg=darkblue}
\setbeamercolor{block title}{fg=white, bg=darkblue}
\setbeamercolor{block body}{fg=black, bg=lightgray!10}

% Убрать навигационные символы
\setbeamertemplate{navigation symbols}{}

% Добавить номера страниц
\setbeamertemplate{footline}{%
    \hfill%
    \usebeamerfont{footline}%
    \usebeamercolor[fg]{footline}%
    \insertframenumber/\inserttotalframenumber%
    \hspace{1em}\vspace{0.5em}%
}

% Установить шрифт
\usepackage[T1]{fontenc}
\usepackage{helvet}
\renewcommand{\familydefault}{\sfdefault}

% Параметры листинга кода
\lstset{
    basicstyle=\ttfamily\small,
    breaklines=true,
    backgroundcolor=\color{white},
    frame=single,
    rulecolor=\color{gray},
    numbers=none,
    keywordstyle=\color{blue},
    commentstyle=\color{gray},
    stringstyle=\color{red},
}

% ЗАМЕНИТЕ ЗДЕСЬ ВАШИ ДАННЫЕ:
\title[Система онбординга]{Система онбординга для новых сотрудников}
\author{
    Фирана Мамедова \\
    \vspace{0.3cm}
    Группа: БПИ-XX \\
    \vspace{0.3cm}
    Руководитель: [ФИО руководителя, должность] \\
    \vspace{0.3cm}
    {\small Консультант: [ФИО консультанта], [Должность] \\
    [Место работы]}
}
\date{\today}

\begin{document}

% ========== СЛАЙД 1: ТИТУЛЬНЫЙ ЛИСТ ==========
\begin{frame}[plain]
    \titlepage
    \addtocounter{framenumber}{-1}
\end{frame}

% ========== СЛАЙД 2: АКТУАЛЬНОСТЬ ==========
\begin{frame}
    \frametitle{Актуальность}
    
    \textbf{Боль пользователя:}
    \begin{itemize}
        \item Первые дни сотрудника — хаотичны и неструктурированы
        
        \item Информация разбросана по разным каналам
        
        \item Сложно понять ожидания и требования компании
        
        \item Нет отслеживания прогресса адаптации
    \end{itemize}
    
    \bigskip
    
    \textbf{Необходимость решения:}
    
    Автоматизированная система онбординга для структурированной адаптации, прозрачного отслеживания прогресса и вовлечения новых сотрудников
\end{frame}

% ========== СЛАЙД 3: ЦЕЛЬ И ЗАДАЧИ ==========
\begin{frame}
    \frametitle{Цель и основные задачи}
    
    \textbf{Цель:} Разработать веб-приложение для автоматизации и оптимизации процесса онбординга новых сотрудников
    
    \bigskip
    
    \textbf{Задачи:}
    \begin{enumerate}
        \item Анализ предметной области и существующих решений
        \item Формирование функциональных и нефункциональных требований
        \item Разработка технического задания на создание системы
        \item Проектирование архитектуры, структуры данных и сценариев
        \item Реализация клиентской и серверной частей программного продукта
        \item Тестирование и оценка соответствия поставленным требованиям
        \item Подготовка системы для внедрения в организации
        \item Анализ результатов и направления дальнейшего развития
    \end{enumerate}
\end{frame}

% ========== СЛАЙД 4: ЗАДАЧА 1 — АНАЛИЗ ПРЕДМЕТНОЙ ОБЛАСТИ ==========
\begin{frame}
    \frametitle{Задача 1: Анализ предметной области}
    
    \textbf{Проведённый анализ:}
    \begin{itemize}
        \item Исследование современного состояния онбординга в компаниях
        \item Анализ существующих платформ и решений (LMS, HR-системы)
        \item Выявление болевых точек и требований пользователей
        \item Определение лучших практик адаптации сотрудников
        \item Исследование технологических стеков для реализации
    \end{itemize}
    
    \bigskip
    
    \textbf{Выводы:}
    \begin{itemize}
        \item Отсутствие интегрированного решения на рынке
        \item Необходимость комплексного подхода к онбордингу
        \item Требование масштабируемости и гибкости системы
    \end{itemize}
\end{frame}

% ========== СЛАЙД 5: ЗАДАЧА 2 — ТРЕБОВАНИЯ ==========
\begin{frame}
    \frametitle{Задача 2: Функциональные и нефункциональные требования}
    
    \textbf{Функциональные требования:}
    \begin{itemize}
        \item Управление пользователями с ролевой системой доступа
        \item Создание и управление модулями обучения
        \item Отслеживание прогресса и аналитика
        \item Real-time уведомления и система коммуникации
        \item Гамификация и система менторства
    \end{itemize}
    
    \bigskip
    
    \textbf{Нефункциональные требования:}
    \begin{itemize}
        \item Масштабируемость и производительность
        \item Безопасность данных и аутентификация
        \item Удобство использования и адаптивность
    \end{itemize}
\end{frame}

% ========== СЛАЙД 6: ЗАДАЧА 3 — ТЕХНИЧЕСКОЕ ЗАДАНИЕ ==========
\begin{frame}
    \frametitle{Задача 3: Разработка технического задания}
    
    \textbf{Определены:}
    \begin{itemize}
        \item Технологический стек: ASP.NET Core, Vue 3, SQL Server
        \item Архитектурный подход: N-tier с разделением на слои
        \item Методология разработки: Agile с итерациями
        \item Сроки и этапы реализации
        \item Критерии приёмки и метрики качества
    \end{itemize}
    
    \bigskip
    
    \textbf{Спецификация:}
    \begin{itemize}
        \item Описание функциональности каждого модуля
        \item API контракты и структуры данных
        \item Требования к производительности и безопасности
    \end{itemize}
\end{frame}

% ========== СЛАЙД 7: ЗАДАЧА 4 — АРХИТЕКТУРА И ПРОЕКТИРОВАНИЕ ==========
\begin{frame}
    \frametitle{Задача 4: Проектирование архитектуры системы}
    
    \textbf{Компоненты системы:}
    \begin{itemize}
        \item Backend: Controllers → Services → Repository → Database
        \item Frontend: Views → Components → Stores → API
        \item Real-time: SignalR Hubs для WebSocket соединений
        \item Инфраструктура: Docker контейнеризация
    \end{itemize}
    
    \bigskip
    
    \textbf{Структуры данных:}
    \begin{itemize}
        \item 14+ сущностей в БД (User, Module, Progress, Notification и др.)
        \item DTO для безопасной передачи данных
        \item Связи один-ко-многим и многие-ко-многим
    \end{itemize}
\end{frame}

% ========== СЛАЙД 8: ЗАДАЧА 5 — РЕАЛИЗАЦИЯ ==========
\begin{frame}
    \frametitle{Задача 5: Реализация клиентской и серверной частей}
    
    \textbf{Backend (ASP.NET Core):}
    \begin{itemize}
        \item JWT-based аутентификация с BCrypt хешированием
        \item 18+ контроллеров API для основной функциональности
        \item SignalR для real-time уведомлений
        \item Entity Framework Core с миграциями БД
    \end{itemize}
    
    \bigskip
    
    \textbf{Frontend (Vue 3 + Vite):}
    \begin{itemize}
        \item 13+ компонентов для интерактивности
        \item Pinia для управления состоянием
        \item Tailwind CSS для стилизации
        \item Интеграция с backend через Axios
    \end{itemize}
\end{frame}

% ========== СЛАЙД 9: ЗАДАЧА 6 — ТЕСТИРОВАНИЕ ==========
\begin{frame}
    \frametitle{Задача 6: Тестирование и оценка соответствия}
    
    \textbf{Проведённое тестирование:}
    \begin{itemize}
        \item Unit-тесты для сервисного слоя
        \item Integration-тесты для API контроллеров
        \item Frontend функциональное тестирование
        \item Тестирование безопасности (SQL-injection, XSS)
    \end{itemize}
    
    \bigskip
    
    \textbf{Результаты:}
    \begin{itemize}
        \item Все функциональные требования выполнены
        \item Система стабильна при нагрузке
        \item Время ответа API < 200ms
        \item Уровень покрытия тестами > 80\%
    \end{itemize}
\end{frame}

% ========== СЛАЙД 10: ЗАДАЧА 7 — ВНЕДРЕНИЕ ==========
\begin{frame}
    \frametitle{Задача 7: Подготовка к внедрению}
    
    \textbf{Выполненные мероприятия:}
    \begin{itemize}
        \item Docker образы и docker-compose для развёртывания
        \item Документация API (Swagger/OpenAPI)
        \item Инструкции администратора и пользователя
        \item Миграция и инициализация БД
        \item Настройка окружений (Dev, Test, Prod)
    \end{itemize}
    
    \bigskip
    
    \textbf{Готовность к внедрению:}
    \begin{itemize}
        \item Система готова к использованию в реальных условиях
        \item Поддержка масштабирования при росте пользователей
        \item Резервное копирование и восстановление данных
    \end{itemize}
\end{frame}

% ========== СЛАЙД 11: ЗАДАЧА 8 — АНАЛИЗ РЕЗУЛЬТАТОВ ==========
\begin{frame}
    \frametitle{Задача 8: Анализ результатов и перспективы}
    
    \textbf{Достигнутые результаты:}
    \begin{itemize}
        \item Полнофункциональная система онбординга реализована
        \item Пользовательский интерфейс интуитивен и удобен
        \item Система демонстрирует высокую производительность
        \item Положительная обратная связь от тестировщиков
    \end{itemize}
    
    \bigskip
    
    \textbf{Направления развития:}
    \begin{itemize}
        \item Интеграция с HRIS системами через API
        \item Расширение AI функционала (анализ прогресса)
        \item Мобильное приложение для iOS/Android
        \item Интеграция видео-контента и вебинаров
    \end{itemize}
\end{frame}

% ========== СЛАЙД 12: ВИДЕО-ДЕМОНСТРАЦИЯ ==========
\begin{frame}
    \frametitle{Демонстрация функциональности}
    
    \centering
    
    \Large \textbf{Видео демонстрация} \\[2em]
    
    {\small \textit{(2-3 минуты)}} \\[1em]
    
    \normalsize
    
    Основные возможности приложения:
    \begin{itemize}
        \item Просмотр модулей обучения
        \item Прохождение тестов и получение результатов
        \item Отслеживание прогресса и аналитика
        \item Real-time уведомления и взаимодействие
        \item Лидерборд и достижения пользователя
    \end{itemize}
\end{frame}

% ========== СЛАЙД 13: ЗАКЛЮЧЕНИЕ ==========
\begin{frame}
    \frametitle{Заключение}
    
    \textbf{Выполненные задачи (совершенная форма):}
    
    \begin{enumerate}[label=\arabic*.]
        \item Проведён анализ предметной области
        
        \item Сформированы функциональные и нефункциональные требования
        
        \item Разработано техническое задание на создание системы
        
        \item Спроектирована архитектура и структура данных
        
        \item Реализованы клиентская и серверная части
        
        \item Проведено тестирование и оценка соответствия требованиям
        
        \item Система подготовлена к внедрению в организации
        
        \item Проведён анализ результатов и определены направления развития
    \end{enumerate}
\end{frame}

% ========== СЛАЙД 14: СТРУКТУРА ВКР ==========
\begin{frame}
    \frametitle{Структура ВКР}
    
    \small
    
    \begin{enumerate}
        \item Введение (стр. \_\_)
        \item Анализ предметной области и требования (стр. \_\_)
        \item Проектирование архитектуры системы (стр. \_\_)
        \item Реализация Backend (стр. \_\_)
        \item Реализация Frontend (стр. \_\_)
        \item Тестирование и развертывание (стр. \_\_)
        \item Заключение (стр. \_\_)
        \item Приложения (стр. \_\_)
    \end{enumerate}
    
    \bigskip
    
    \small \textit{Номера страниц добавить после завершения ВКР}
\end{frame}

\end{document}